\preliminar{Resumen}

\noindent
El presente proyecto aborda la ineficiencia del proceso de venta y reserva
de paquetes turísticos en una agencia de turismo de la ciudad de Potosí,
donde la verificación de cupos y la actualización de la información no se
realizan de manera uniforme entre el personal. El objetivo general es
desarrollar un sistema web que permita administrar paquetes, fechas de
salida y reservas, validar la disponibilidad de cupos y mejorar la
confirmación de viajes. La propuesta se sustenta en los antecedentes de los
sistemas de reservas, en el marco de la gestión de paquetes turísticos y en
criterios de usabilidad y seguridad para aplicaciones web. La metodología
prevé entrevistas, observación y revisión autorizada de documentos en la
agencia seleccionada, con el fin de diagnosticar el proceso actual y
diseñar la solución de acuerdo con sus necesidades. Se espera que el
sistema contribuya a reducir inconsistencias en la disponibilidad,
disminuir el tiempo de confirmación y ofrecer información confiable para la
toma de decisiones.

\vspace{0.5\baselineskip}
\noindent\textbf{Palabras clave:} agencias de viaje, paquetes turísticos,
reservas, sistemas web, control de disponibilidad
